\documentclass[11pt]{article}
\usepackage{amsmath,amssymb,amsthm}
\usepackage{algorithm}
\usepackage[noend]{algpseudocode}
\usepackage{geometry}
\geometry{margin=1in}
\newtheorem{theorem}{Theorem}
\newtheorem{lemma}{Lemma}
\newtheorem{assumption}{Assumption}
\newcommand{\E}{\mathbb{E}}
\newcommand{\R}{\mathbb{R}}

\begin{document}

\section*{Algorithm Name}
\textbf{Stochastic Gradient Descent with Fixed Step Size (SGD‑FS)}  

\section*{Mathematical Formulation}
Let $f:\R^{d}\rightarrow\R$ be the population objective
\[
f(w)=\mathbb{E}_{\xi\sim\mathcal{D}}\!\big[\,\ell(w;\xi)\,\big],
\]
where $\ell(w;\xi)$ is the loss incurred on a random data point $\xi$.  
Given an initial point $w_{0}\in\R^{d}$ and a constant stepsize $\eta>0$, the iterates are defined for $t=0,1,2,\dots$ by
\begin{align}
\label{eq:update}
w_{t+1}&=w_{t}-\eta\,g_{t},\qquad 
g_{t}\;:=\;\nabla\ell\big(w_{t};\xi_{t}\big),
\end{align}
where $\xi_{t}$ is drawn i.i.d. from $\mathcal{D}$ (or, equivalently, a minibatch is sampled and $g_{t}$ is the average gradient over that minibatch).  

\section*{Pseudocode}
\begin{algorithm}[H]
\caption{SGD‑FS}
\begin{algorithmic}[1]
\State \textbf{Input:} stepsize $\eta>0$, initial point $w_{0}\in\R^{d}$, maximum iterations $T$.
\For{$t=0$ \textbf{to} $T-1$}
    \State Sample a minibatch $\mathcal{B}_{t}$ (or a single data point) from the training distribution.
    \State Compute stochastic gradient $g_{t}= \frac{1}{|\mathcal{B}_{t}|}\sum_{\xi\in\mathcal{B}_{t}}\nabla\ell(w_{t};\xi)$.
    \State Update $w_{t+1}=w_{t}-\eta\,g_{t}$.
\EndFor
\State \textbf{Output:} $w_{T}$ (or optionally the average $\bar w_T=\frac1T\sum_{t=1}^{T}w_t$).
\end{algorithmic}
\end{algorithm}

\section*{Dry Run Example}
Consider the one–dimensional quadratic
\[
f(w)=(w-3)^{2},\qquad \nabla f(w)=2(w-3).
\]
Assume the stochastic gradient is exact (i.e. $g_{t}=\nabla f(w_{t})$) and choose
\[
w_{0}=0,\qquad \eta=0.1.
\]

\begin{center}
\begin{tabular}{c|c|c|c}
$t$ & $w_{t}$ & $g_{t}=2(w_{t}-3)$ & $f(w_{t})$\\ \hline
0 & $0$ & $-6$ & $9$\\
1 & $w_{1}=0-0.1(-6)=0.6$ & $2(0.6-3)=-4.8$ & $(0.6-3)^{2}=5.76$\\
2 & $w_{2}=0.6-0.1(-4.8)=1.08$ & $2(1.08-3)=-3.84$ & $(1.08-3)^{2}=3.6864$\\
3 & $w_{3}=1.08-0.1(-3.84)=1.464$ & $2(1.464-3)=-3.072$ & $(1.464-3)^{2}=2.3650$
\end{tabular}
\end{center}
The objective value decreases at each iteration, illustrating the contraction property of the method on a strongly convex quadratic.

\newpage
\section*{Assumptions}
\begin{assumption}[Strong Convexity]\label{ass:strong}
The population objective $f$ is $\mu$–strongly convex, i.e. for all $u,v\in\R^{d}$
\[
f(v)\ge f(u)+\langle\nabla f(u),v-u\rangle+\frac{\mu}{2}\|v-u\|^{2},
\]
with a constant $\mu>0$.
\end{assumption}

\begin{assumption}[Smoothness]\label{ass:smooth}
$f$ has $L$–Lipschitz gradients (i.e. $L$–smooth):
\[
\|\nabla f(u)-\nabla f(v)\|\le L\|u-v\|,\qquad\forall u,v\in\R^{d},
\]
for some $L\ge\mu$.
\end{assumption}

\begin{assumption}[Unbiased Stochastic Gradient]\label{ass:unbiased}
For every $t$, conditional on $w_{t}$ the stochastic gradient satisfies
\[
\E\!\left[g_{t}\mid w_{t}\right]=\nabla f(w_{t}).
\]
\end{assumption}

\begin{assumption}[Bounded Variance]\label{ass:variance}
There exists $\sigma^{2}\ge0$ such that
\[
\E\!\left[\|g_{t}-\nabla f(w_{t})\|^{2}\mid w_{t}\right]\le\sigma^{2},
\qquad\forall t.
\]
\end{assumption}

\begin{assumption}[Stepsize]\label{ass:stepsize}
The constant stepsize fulfills
\[
0<\eta<\frac{2}{\mu+L}.
\]
\end{assumption}

\section*{Main Convergence Theorem}
\begin{theorem}[Linear convergence to a neighbourhood]\label{thm:linear}
Under Assumptions \ref{ass:strong}–\ref{ass:stepsize}, the iterates generated by \eqref{eq:update} satisfy for all $t\ge0$
\[
\boxed{\;
\E\!\big[\|w_{t}-w^{*}\|^{2}\big]\;\le\;\big(1-\eta\mu\big)^{t}\|w_{0}-w^{*}\|^{2}
\;+\;\frac{\eta\sigma^{2}}{\mu}\big(1-(1-\eta\mu)^{t}\big)\;},
\tag{1}
\]
where $w^{*}=\arg\min f$.
Consequently,
\[
\limsup_{t\to\infty}\E\!\big[\|w_{t}-w^{*}\|^{2}\big]\le\frac{\eta\sigma^{2}}{\mu},
\]
i.e. the method converges geometrically to a ball of radius $O(\sqrt{\eta})$ around the optimum.
\end{theorem}

\section*{Theoretical Properties}
\begin{itemize}
\item \textbf{Stability:} The contraction factor $1-\eta\mu$ is strictly less than one because of \eqref{ass:stepsize}. Hence the recursion is stable.
\item \textbf{Hyperparameter Sensitivity:} Larger $\eta$ speeds up the contraction $1-\eta\mu$ but also inflates the steady‑state error term $\frac{\eta\sigma^{2}}{\mu}$. The optimal trade‑off (minimising the RHS of \eqref{thm:linear}) is $\eta^{\star}= \frac{2}{\mu+L}$, the largest admissible stepsize.
\item \textbf{Comparison to Deterministic Gradient Descent:} If $\sigma^{2}=0$ (exact gradients), \eqref{thm:linear} reduces to the classic linear rate $\|w_{t}-w^{*}\|^{2}\le(1-\eta\mu)^{t}\|w_{0}-w^{*}\|^{2}$.
\item \textbf{Effect of Minibatch Size:} Increasing the minibatch reduces $\sigma^{2}$ proportionally, shrinking the neighbourhood of convergence.
\end{itemize}

\section*{Discussion}
The theorem shows that for strongly convex and smooth objectives, a fixed stepsize SGD enjoys a geometric decay of the expected squared error up to a variance‑dependent floor. This matches the textbook result for constant‑stepsize SGD and justifies the widespread practice of using a modest constant learning rate when the data are noisy but the problem is well‑conditioned.

\newpage
\section*{Preliminaries (Definitions and Lemmas)}
\begin{lemma}[Three‑point identity for smooth functions]\label{lem:smooth}
If $f$ is $L$‑smooth then for any $u,v\in\R^{d}$
\[
\|\nabla f(u)-\nabla f(v)\|^{2}\le L\langle\nabla f(u)-\nabla f(v),\,u-v\rangle.
\tag{2}
\]
\end{lemma}
\begin{proof}
Directly from the $L$‑smoothness inequality
$f(v)\le f(u)+\langle\nabla f(u),v-u\rangle+\frac{L}{2}\|v-u\|^{2}$
and the analogous inequality with $u$ and $v$ swapped; subtracting yields the claimed bound. \qedhere
\end{proof}

\begin{lemma}[Quadratic bound from strong convexity]\label{lem:strong}
For a $\mu$‑strongly convex $f$,
\[
\langle\nabla f(u)-\nabla f(v),\,u-v\rangle\ge\mu\|u-v\|^{2},
\qquad\forall u,v\in\R^{d}.
\tag{3}
\]
\end{lemma}
\begin{proof}
Apply the definition of strong convexity at $u$ and $v$, then add the two inequalities and rearrange. \qedhere
\end{proof}

\section*{Main Proof (Step‑by‑Step Derivation)}
We prove Theorem \ref{thm:linear}.  
All expectations are taken with respect to the randomness up to iteration $t$ unless otherwise noted.

\bigskip
\noindent\textbf{Step 1.} Start from the update \eqref{eq:update} and expand the squared distance to the optimum $w^{*}$:
\[
\|w_{t+1}-w^{*}\|^{2}
= \|w_{t}-\eta g_{t}-w^{*}\|^{2}
= \|w_{t}-w^{*}\|^{2} -2\eta\langle g_{t},w_{t}-w^{*}\rangle +\eta^{2}\|g_{t}\|^{2}.
\tag{4}
\]

\noindent\textbf{Step 2.} Take conditional expectation given $w_{t}$ and use unbiasedness (Assumption \ref{ass:unbiased}):
\[
\E\!\big[\|w_{t+1}-w^{*}\|^{2}\mid w_{t}\big]
= \|w_{t}-w^{*}\|^{2} -2\eta\langle \nabla f(w_{t}),w_{t}-w^{*}\rangle
+ \eta^{2}\E\!\big[\|g_{t}\|^{2}\mid w_{t}\big].
\tag{5}
\]

\noindent\textbf{Step 3.} Decompose $\|g_{t}\|^{2}$ into variance and bias parts:
\[
\|g_{t}\|^{2}= \|\nabla f(w_{t})\|^{2}
+ \|g_{t}-\nabla f(w_{t})\|^{2}
+2\langle\nabla f(w_{t}),g_{t}-\nabla f(w_{t})\rangle.
\]
Taking conditional expectation, the cross term vanishes (by unbiasedness) and the variance term is bounded by $\sigma^{2}$ (Assumption \ref{ass:variance}). Hence
\[
\E\!\big[\|g_{t}\|^{2}\mid w_{t}\big]\le \|\nabla f(w_{t})\|^{2}+\sigma^{2}.
\tag{6}
\]

\noindent\textbf{Step 4.} Substitute \eqref{6} into \eqref{5}:
\[
\E\!\big[\|w_{t+1}-w^{*}\|^{2}\mid w_{t}\big]
\le \|w_{t}-w^{*}\|^{2}
-2\eta\langle \nabla f(w_{t}),w_{t}-w^{*}\rangle
+\eta^{2}\|\nabla f(w_{t})\|^{2}
+\eta^{2}\sigma^{2}.
\tag{7}
\]

\noindent\textbf{Step 5.} Relate the inner product to the squared norm using strong convexity (Lemma \ref{lem:strong}) and smoothness (Lemma \ref{lem:smooth}). First, from Lemma \ref{lem:strong}
\[
\langle \nabla f(w_{t})- \nabla f(w^{*}),\,w_{t}-w^{*}\rangle\ge\mu\|w_{t}-w^{*}\|^{2}.
\tag{8}
\]
Since $\nabla f(w^{*})=0$, (8) simplifies to
\[
\langle \nabla f(w_{t}),\,w_{t}-w^{*}\rangle\ge\mu\|w_{t}-w^{*}\|^{2}.
\tag{9}
\]

\noindent\textbf{Step 6.} Upper‑bound $\|\nabla f(w_{t})\|^{2}$ via smoothness. By Lemma \ref{lem:smooth} and (8):
\[
\|\nabla f(w_{t})\|^{2}
= \|\nabla f(w_{t})-\nabla f(w^{*})\|^{2}
\le L\langle \nabla f(w_{t})-\nabla f(w^{*}),\,w_{t}-w^{*}\rangle
\le L\langle \nabla f(w_{t}),\,w_{t}-w^{*}\rangle .
\tag{10}
\]

\noindent\textbf{Step 7.} Plug (9) and (10) into (7):
\[
\begin{aligned}
\E\!\big[\|w_{t+1}-w^{*}\|^{2}\mid w_{t}\big]
&\le \|w_{t}-w^{*}\|^{2}
-2\eta\mu\|w_{t}-w^{*}\|^{2}
+ \eta^{2}L\langle \nabla f(w_{t}),w_{t}-w^{*}\rangle
+\eta^{2}\sigma^{2}\\[4pt]
&\le \bigl(1-2\eta\mu\bigr)\|w_{t}-w^{*}\|^{2}
+ \eta^{2}L\bigl(\tfrac{1}{\mu}\bigr)\mu\|w_{t}-w^{*}\|^{2}
+\eta^{2}\sigma^{2}
\quad\text{(using (9) to replace the inner product)}\\[4pt]
&= \bigl(1-2\eta\mu + \eta^{2}L\mu^{-1}\mu\bigr)\|w_{t}-w^{*}\|^{2}
+\eta^{2}\sigma^{2}\\[4pt]
&= \bigl(1-2\eta\mu + \eta^{2}L\bigr)\|w_{t}-w^{*}\|^{2}
+\eta^{2}\sigma^{2}.
\end{aligned}
\tag{11}
\]

\noindent\textbf{Step 8.} Choose $\eta$ such that the coefficient in front of $\|w_{t}-w^{*}\|^{2}$ is strictly smaller than one. Under Assumption \ref{ass:stepsize},
\[
1-2\eta\mu+\eta^{2}L \;=\; 1-\eta\mu\bigl(2-\tfrac{\eta L}{\mu}\bigr)
\;\le\; 1-\eta\mu,
\tag{12}
\]
because $0<\eta<\frac{2}{\mu+L}$ implies $\eta L\le 2-\eta\mu$.  
Thus,
\[
\E\!\big[\|w_{t+1}-w^{*}\|^{2}\mid w_{t}\big]
\le (1-\eta\mu)\|w_{t}-w^{*}\|^{2}+\eta^{2}\sigma^{2}.
\tag{13}
\]

\noindent\textbf{Step 9.} Take full expectation and unfold the recursion. Define $a_{t}:=\E\!\big[\|w_{t}-w^{*}\|^{2}\big]$. From (13):
\[
a_{t+1}\le (1-\eta\mu)a_{t}+\eta^{2}\sigma^{2}.
\tag{14}
\]
Iterating (14) yields a standard linear difference equation solution:
\[
\begin{aligned}
a_{t}
&\le (1-\eta\mu)^{t}a_{0}
+ \eta^{2}\sigma^{2}\sum_{k=0}^{t-1}(1-\eta\mu)^{k}\\[4pt]
&= (1-\eta\mu)^{t}\|w_{0}-w^{*}\|^{2}
+ \eta^{2}\sigma^{2}\,\frac{1-(1-\eta\mu)^{t}}{\eta\mu}\\[4pt]
&= (1-\eta\mu)^{t}\|w_{0}-w^{*}\|^{2}
+ \frac{\eta\sigma^{2}}{\mu}\Bigl(1-(1-\eta\mu)^{t}\Bigr).
\end{aligned}
\tag{15}
\]

\noindent\textbf{Step 10.} Equation (15) is precisely the bound (1) claimed in Theorem \ref{thm:linear}. Taking $t\to\infty$ gives the steady‑state error term $\frac{\eta\sigma^{2}}{\mu}$ and completes the proof. \qed

\section*{Rate Derivation (With Constants)}
From (15) the dominant term for large $t$ is the geometric factor $(1-\eta\mu)^{t}\approx e^{-\eta\mu t}$. Hence the expected distance decays at a linear (geometric) rate with contraction constant $q:=1-\eta\mu\in(0,1)$. The asymptotic neighbourhood radius is
\[
R:=\sqrt{\frac{\eta\sigma^{2}}{\mu}}.
\]
If the stochastic gradient is exact ($\sigma^{2}=0$) the neighbourhood collapses to the optimum and the bound reduces to
\[
\E\!\big[\|w_{t}-w^{*}\|^{2}\big]\le (1-\eta\mu)^{t}\|w_{0}-w^{*}\|^{2}.
\]

\section*{Special Cases}
\begin{itemize}
\item \textbf{Deterministic Gradient Descent (full batch).} $\sigma^{2}=0$ $\Rightarrow$ linear convergence to the exact solution.
\item \textbf{Mini‑batch size $B$.} If the variance of a single sample is $\sigma_{1}^{2}$, then $\sigma^{2}=\sigma_{1}^{2}/B$, so the steady‑state error shrinks as $1/\sqrt{B}$.
\item \textbf{Optimal constant stepsize.} Maximising the contraction $1-\eta\mu$ while keeping $\eta$ admissible gives $\eta^{\star}=2/(\mu+L)$. Plugging $\eta^{\star}$ into (1) yields
\[
\E\!\big[\|w_{t}-w^{*}\|^{2}\big]\le
\Bigl(\frac{L-\mu}{L+\mu}\Bigr)^{t}\|w_{0}-w^{*}\|^{2}
+ \frac{2\sigma^{2}}{\mu(L+\mu)}.
\]
\end{itemize}

\end{document}